\begin{table}[!htbp]\centering
\caption{Entrada y permanencia en la formalidad en el panel anual}\label{tab:journal_panel}
\begin{threeparttable}\small
\setlength{\tabcolsep}{5pt}
\begin{tabular}{lrrrr}
\toprule
Ventana & DiD (pp) & IC 95\% & $p$ & $N$ \\
\midrule
\multicolumn{5}{l}{\emph{Entrada a empleo privado formal}} \\
Cruza la reforma & -2.54 & [-5.14, 0.06] & 0.055 & 33076 \\
Posterior a la reforma & -4.66 & [-8.63, -0.68] & 0.024 & 33076 \\
\multicolumn{5}{l}{\emph{Permanencia en empleo privado formal}} \\
Cruza la reforma & 8.89 & [1.14, 16.64] & 0.026 & 7011 \\
Posterior a la reforma & -1.41 & [-5.40, 2.58] & 0.472 & 7011 \\
\bottomrule
\end{tabular}
\begin{tablenotes}\footnotesize
\item ENAHO Panel 2020--2024. Coeficientes de baja educación por ventana de la ecuación~\eqref{eq:journal_panel}; referencia: pares previos. Entrada: origen asalariado informal o independiente no clasificado como formal; permanencia: origen privado formal. Factores de expansión del panel, controles del año inicial, efectos fijos de par de años y departamento e indicadores comunes de ventana. IC con errores agrupados por departamento. N cuenta pares-persona, no individuos únicos. Tasas previas de baja educación: 4.51\% de entrada y 62.64\% de permanencia.
\end{tablenotes}\end{threeparttable}\end{table}
